\documentclass[11pt,a4paper]{article}
\usepackage[margin=1in]{geometry}
\usepackage{setspace}
\usepackage{parskip}
\usepackage{enumitem}
\setlist[itemize]{leftmargin=1.5em}
\onehalfspacing

\title{Speaker Script\\Earnings Surprise, Firm Valuation, and Market Efficiency: Evidence from China}
\author{FIN803 Corporate Finance Team}
\date{}

\begin{document}
\maketitle

\section*{Slide 1: Title}
Good afternoon everyone. We are presenting our FIN803 Corporate Finance project on earnings surprise, valuation adjustment, and market efficiency in China A-share stocks.

\section*{Slide 2: Why This Matters in Corporate Finance}
In corporate finance, earnings disclosures matter because they update beliefs about future cash flows, profitability, and firm value. If markets are reasonably efficient, stock prices should respond when these disclosures contain new information. So this project asks a classic finance question: how quickly and how clearly does disclosure news enter prices?

\section*{Slide 3: Research Question and Prediction}
Our research question is straightforward. Do firms with more positive earnings surprise experience better short-run abnormal returns around the disclosure date? The prediction from standard valuation logic is yes. If the disclosed earnings news is better than expected, then prices should adjust upward. At the same time, we want to be careful, because a clean test depends heavily on how expectations are measured.

\section*{Slide 4: Sample and Empirical Setup}
We focus our headline analysis on earnings preannouncements. The reason is that preannouncements are the cleanest event type for identifying the information shock. They are easier to interpret than mixing together revisions, express releases, and formal releases, which can reflect different information environments. We then study short-run abnormal returns over several windows after the event, from 3 to 20 trading days.

\section*{Slide 5: How We Measure Earnings Surprise}
The main methodological challenge is measuring surprise in a credible way. We benchmark the disclosure against analyst expectations that are matched as closely as possible to the same reporting period. We also compare three versions of surprise: level surprise, percentage surprise, and standardized surprise. This comparison matters because the empirical conclusion can change depending on how the surprise variable is scaled, so we do not want to rely on only one definition.

\section*{Slide 6: Main Result: Directional but Modest Evidence}
This slide shows the main visual result. In the cleaner baseline sample, firms with higher earnings surprise tend to outperform firms with lower earnings surprise after the disclosure. That is directionally consistent with short-run valuation adjustment. But the important point is that the separation is modest rather than dramatic. So the figure supports the theory, but it does not justify a strong headline claim on its own.

\section*{Slide 7: Results Depend on Specification Choices}
The second result is that the estimated relation is sensitive to specification choices. When we change how surprise is scaled, or how tightly analyst expectations are matched to the event, the estimated coefficient can move noticeably. This is important because it means the result is not fully robust across all reasonable designs. In other words, the evidence is informative, but it is not strong enough for us to overstate the conclusion.

\section*{Slide 8: What We Learn About Valuation and Efficiency}
What do we learn from this? First, the cleaner baseline is broadly consistent with the idea that earnings disclosures affect valuation in the short run. Second, the evidence is not strong enough to claim a large or universal market inefficiency story. A more defensible interpretation is that there is some short-window price adjustment, but the measured strength depends heavily on expectation construction and empirical design.

\section*{Slide 9: Why We Keep the Conclusion Cautious}
This is why our final interpretation is deliberately cautious. The headline result is weak rather than decisive. The inference changes when we alter surprise measurement, matching quality, or event composition. That means empirical discipline matters as much as statistical significance here. For a class project, we think it is better to present an honest result with clear logic than to force a stronger claim than the evidence can support.

\section*{Slide 10: Conclusion}
To conclude, we take away three main points. First, earnings-related disclosures do appear to matter for valuation in the cleaner preannouncement sample. Second, the evidence is directionally supportive, but modest and sensitive to measurement choices. Third, the most credible interpretation is therefore cautious: this is suggestive evidence on disclosure and market efficiency, not a definitive headline result. Thank you.

\end{document}
